%!TEX root = ./main.tex
\chapter{绪论}
\label{chap:introduction}

\section{研究背景}

今天，计算机技术和互联网络的飞速发展把社会的信息化进程推向了一个全新的阶段，信息的传递与交流已经成为整个现代社会生活运作的重要基础，电子可读文本大量涌现并成为网络时代主要的信息载体和人们的生活中不可缺的一部分。随着信息化进代的来临，自然语言处理技术已逐渐成为一项大从的迫切需求，计算语言学的研究也越来越受到人们的重视。

自然语言分析技术（Natural Language Parsing）一直是计算语言学领域一个基础性的研究课题。大部分自然语言处理系统，包括机器翻译，文本理解，信息的检索与过滤，语音识别与合成，都毫无疑问地会从高质量的分析技术中受益。从科学的观点来看，计算机的自然语言分析过程是对人类语言理解过程的模拟：即根据一定的语言知识，通常是一具由规则、树或图组成的形式文法系统，将输入句子的一维线性结构赋予某种二维平面结构解释；从人工工智能研究的角度来讲，这是一个基于推理的问题求解过程，分析方法则对应了其推理控制策略。~\cite{guo_2023_research,thakkar_2018_performance,zheng_2019_nutbaas}

Testing测试12345章节~\upcite{guo_2023_research}

\subsection{结构歧义}

然而与形式语言的分析相比，应用计算机进行自然语言的结构分析却远非一件容易的工作。

\subsubsection{测试AA}

然而与形式语言的分析相比，应用计算机进行自然语言的结构分析却远非一件容易的工作。
